
\chapter{Product Specifications}
In last session's tutorial, we translated customer statements into customer needs. In this session, we will continue working with these needs, aiming to transform them into measurable metrics that the Product Development (PD) team can utilize effectively.

\subsection{In-class Activities}

\subsubsection{Activity 1: Create a needs and importance Table}
First, please prepare a table with columns labelled Number, Sub-component, Needs, and Importance. Below is an example of how the table might look as shown in Table \ref{tab:suspension_needs}:

\begin{table}[h]
    \centering
    \caption{Suspension system needs and importance}
    \label{tab:suspension_needs}
    \begin{tabular}{|c|c|c|c|}
        \hline
        No & Sub-component & Need & Imp. \\ \hline
        1 & The suspension & Reduces vibration to the hands. & 3 \\ \hline
        2 & The suspension & Allows easy traversal of slow, difficult terrain. & 2 \\ \hline
    \end{tabular}
\end{table}

\noindent
You should list the sub-components that your group has identified from analyzing customer needs. Also, for each need you identify, assign an importance level on a scale from 1 to 5, where 1 means less important and 5 means very important.

\subsubsection{Activity 2: Defining Metrics for Smooth Product Functionality} % if available
Now, let's find the appropriate metrics for each of the customer needs. Due to time constraints, we will use the LLM (language learning model) to assist us. Remember, this course is not focused on prompt engineering or LLM, but we are using it here to save time. You might consider using the following prompt: \href{https://gist.github.com/balandongiv/b63414ea1f72e311bfdd2f1bdc5468dd}{click me}.


The LLM may produce different results depending on the trial and which service (ChatGPT, Gemini, Claude, etc) you used. It is important for you to assess the accuracy and appropriateness of these results. Then, map the output in the Table \ref{table:bicycle_need_metric_unit_justification} according to your evaluation.



\begin{landscape}
\begin{footnotesize}
\begin{longtable}{|c|p{4cm}|p{5cm}|c|p{6.5cm}|}
\caption{Bicycle Suspension Needs, Metrics, Units, and Justification} \label{table:bicycle_need_metric_unit_justification} \\
\hline
\textbf{No} & \textbf{Need} & \textbf{Metric} & \textbf{Units} & \textbf{Justification} \\
\hline
\endfirsthead

\multicolumn{5}{c}%
{\tablename\ \thetable\ -- \textit{Continued from previous page}} \\
\hline
\textbf{No} & \textbf{Need} & \textbf{Metric} & \textbf{Units} & \textbf{Justification} \\
\hline
\endhead

\hline \multicolumn{5}{r}{\textit{Continued on next page}} \\
\endfoot

\hline
\endlastfoot

\multirow{3}{*}{1} 
  & \multirow{3}{4.5cm}{Reduces vibration to the hands} 
    & 1. Attenuation from dropout to handlebar at 10\,Hz & dB & High-frequency vibrations are most noticeable at the hands; reduction improves rider comfort. \\
\cline{3-5}
  & & 2. Maximum value from the Monster & g & Captures peak impact events to evaluate suspension performance. \\
\cline{3-5}
  & & 3. Minimum descent time on test track & s & Demonstrates performance under real trail conditions. \\
\hline
2 & Allows easy traversal of slow, difficult terrain 
  & 1. Spring preload & N & Adjustability of preload helps with traction and control on rough, slow trails. \\
\hline
\multirow{3}{*}{3} 
  & \multirow{3}{4.5cm}{Enables high-speed descents on bumpy trails} 
    & 1. Attenuation from dropout to handlebar at 10\,Hz & dB & Reducing vibration at high speeds helps maintain control and comfort. \\
\cline{3-5}
  & & 2. Maximum value from the Monster & g & Measures extreme shock loads during descents. \\
\cline{3-5}
  & & 3. Minimum descent time on test track & s & Indicates efficiency and capability over fast, bumpy terrain. \\
\hline
4 & Allows sensitivity adjustment 
  & 1. Damping coefficient adjustment range & N$\cdot$s/m & Provides tunability for rider preference and terrain type. \\
\hline
5 & Preserves the steering characteristics of the bike 
  & 1. Maximum travel (26-in. wheel) & mm & Excessive suspension travel can affect geometry and handling. \\
\hline
\multirow{2}{*}{6} 
  & \multirow{2}{4.5cm}{Remains rigid during hard cornering} 
    & 1. Spring preload & N & A stiffer spring resists compression and maintains geometry. \\
\cline{3-5}
  & & 2. Lateral stiffness at the tip & kN/m & Prevents lateral flexing which would impair handling under load. \\
\hline
7 & Is lightweight 
  & 1. Total mass & kg & Lower mass improves bike efficiency and handling. \\
\hline
8 & Provides stiff mounting points for the brakes 
  & 1. Lateral stiffness at brake pivots & kN/m & Ensures braking forces do not deform the frame and affect control. \\
\hline

\end{longtable}
\end{footnotesize}
\end{landscape}


Just to make life interesting, you could modify the prompt to generate 'Subjective Metrics for Subjective Needs'.

\subsubsection{Activity 3: Defining Metrics for Smooth Product Functionality} % if available
As I often remind you, do not fully trust everything that the LLM produces. According to the textbook, we should follow certain guidelines when creating product specifications. Therefore, please add a new column to your table and check if any suggestions meet the following criteria:
\FloatBarrier
\begin{enumerate}
    \item Metrics that Reflect Customer Needs.
    \item Dependent Metrics for Design Freedom
    \item Practical Metrics for Product Development
    \item Subjective Metrics for Subjective Needs
    \item Metrics for Market Comparison
\end{enumerate}
\FloatBarrier
Evaluate whether the response from the LLM fully or partially meets any of these guidelines. In the new column, mark the numbers 1, 2, 3, 4, or 5 to indicate which guideline(s) each suggestion aligns with.

\subsubsection{Activity 4: Translate into need-metric matrix (Optional)} % if available
This is an optional task that you can choose to do at home. Next, using Microsoft Excel, you will need to rearrange the table from Activity 3 into a new format called a 'need-metric matrix'. Your matrix should look similar to the example provided in Table \ref{table: need to metric pair}



\input{session8/need_metric_pair}

\subsubsection{Activity 5: Collect Competitive Benchmarking Information} % if available
Next, we will assess the metric values used by various competing companies. You may create a table similar to the following Table \ref{table:metric comparison} with columns representing each competitor's product. For this table, each row should correspond to a metric we identified earlier.



\input{session8/comparison_suspension_metric_bicycle.tex}


% \end{landscape}
Next, we will Input the metric values for each competitor's product in the respective columns. In actual product development process, these values should be extracted from competitors' catalogs and supporting literature. However, due to time constraints, we will once again use the Language Learning Model (LLM) to help us define the appropriate specifications for each customer need. It's important to remember that this class does not focus on prompt engineering or the mechanics of LLMs; however, we utilize this tool to efficiently manage our time.


You might consider using the following prompt to engage the LLM: \href{https://gist.github.com/balandongiv/1c360c0ab5658f611c870ea06b4aabcc}{clickme}.  The LLM is expected to generate outputs in various formats depending on the trial and platform used. You must critically evaluate the accuracy and appropriateness of these outputs using your judgment, then document these metrics in the table as instructed.


\subsubsection{Activity 6: Collect Competitive Benchmarking Information} % if available

For better visualization, we can represent these values using graphs. This helps us to clearly see the range of metric values across different competitive products. 

You may consider the following prompt \href{https://gist.github.com/balandongiv/3f06a6bdbf30c20dbcfd6877a46c8373}{click me} to generate a bar chart to display the values adopted by various competitors for a specific metric. Alternatively, you can generate a grouped bar chart \href{https://gist.github.com/balandongiv/dc9635a4feaf42694f8f5054b10f13cf}{click me} to show the values adopted by various competitors across a range of metrics related to the product.
If you are using the free version of ChatGPT, it might provide you with Python code instead of a direct graph. You can either:



\begin{enumerate}
    \item Ask a friend who has taken a computer programming class to run the code and share the results with you.
    \item Use  \href{https://www.codabrainy.com/en/python-compiler/}{codabrainy} or \href{https://trinket.io/embed/python3}{} (trinke), which does not require any installation, to run the code yourself.
\end{enumerate}

This exercise is intended to help you understand how to visually compare different values used by competing products. Remember, the image created by the LLM might not always be accurate, so it is important to verify the output yourself. You can also use Microsoft Excel to create these graphs for more precise results.

